%%%%%%%%%%%%%%%%%%%%%%%%%%%%%%%%%%%%%%%%%%%%%%%%%%%%%%%%%%%%%%%%%%%%%%%%%%%
%
% Generic template for TFC/TFM/TFG/Tesis
%
% $Id: estudioTeorico.tex,v 1.5 2015/06/05 00:05:19 macias Exp $
%
% By:
%  + Javier Macías-Guarasa.
%    Departamento de Electrónica
%    Universidad de Alcalá
%  + Roberto Barra-Chicote.
%    Departamento de Ingeniería Electrónica
%    Universidad Politécnica de Madrid
% 
% Based on original sources by Roberto Barra, Manuel Ocaña, Jesús Nuevo, Pedro Revenga, Fernando Herránz and Noelia Hernández. Thanks a lot to all of them, and to the many anonymous contributors found (thanks to google) that provided help in setting all this up.
%
% See also the additionalContributors.txt file to check the name of additional contributors to this work.
%
% If you think you can add pieces of relevant/useful examples, improvements, please contact us at (macias@depeca.uah.es)
%
% You can freely use this template and please contribute with comments or suggestions!!!
%
%%%%%%%%%%%%%%%%%%%%%%%%%%%%%%%%%%%%%%%%%%%%%%%%%%%%%%%%%%%%%%%%%%%%%%%%%%%

\chapter{Estado del arte}
\label{cha:estudio-teorico}

Las diferentes resoluciones del Problema del Viajante varían significativamente en función del método 
empleado para ello. En el ámbito de las resoluciones exactas, existen métodos tan sencillos como la 
fuerza bruta, la cual es capaz de encontrar la solución exacta en un espantoso $O(n!)$ \cite{baidoo2016}.
Los algoritmos voraces, los cuales son el primer intento de optimización del problema, presentan una 
complejidad de $O(n^2 \log n)$, aunque no garantizan el óptimo global \cite{abdulkarim2015}. El algoritmo
de ramificación y acotamiento\footnote{Más conocido como \textit{branch and bound}, su término inglés.}
obtiene una complejidad exponencial en la media, aunque puede empeorar a una factorial en el peor de los
casos, lo cual implica, esencialmente, recorrer todo el árbol de búsqueda \cite{thakoor2009}. En cuanto a 
la programación dinámica, se establece el problema mediante $2^n$ posibles subconjuntos de $n$ ciudades o
nodos, debiéndose resolver cada uno de los $n \cdot 2^n$ subproblemas. Cada uno de ellos requiere evaluar,
además, otro subconjunto de $n$ problemas más pequeños, llegando a una complejidad temporal media de 
$O(n^2 \cdot 2^n)$ \cite{bouman2018}. Por último, el algoritmo del vecino más cercano\footnote{
\textit{Nearest Neighbour} en inglés.} presenta una complejidad computacional de $O(n^2)$, debido a que el 
recorrido por todos los nodos se hace en $n$ pasos, pero se debe realizar desde cada uno de los $n$ 
nodos. Este resultado parece contradecir la premisa de que para un algoritmo NP-completo no existe ningún 
algoritmo que lo solucione en tiempo polinómico, pero lo cierto es que este algoritmo no genera un 
resultado óptimo. Se puede ver un caso práctico en \cite{sahalot2014}, donde un conjunto de datos muy 
pequeño\footnote{5 nodos.} basta para afirmar que la fuerza bruta es una mejor técnica de resolución que
el algoritmo del vecino más próximo. Aun así, esta técnica es un paso hacia los métodos más heurísticos y
aproximados que se van a comentar en los siguientes apartados.

En \cite{baidoo2016} se lleva a cabo un estudio experimental con estos cinco métodos de resolución, donde
se verifican las afirmaciones realizadas hasta el momento. Empleando un conjunto de datos aleatorizado, 
y una implementación adaptada para cada algoritmo en Java, se concluyó lo siguiente:
\begin{itemize}
  \item El óptimo global se encuentra mediante fuerza bruta, con la clara desventaja de que es el método 
más ineficiente.
  \item El algoritmo voraz y el algoritmo del vecino más próximo propocionan soluciones muy próximas al 
  óptimo, en un margen de tiempo aceptable.
  \item Ramificación y acotamiento, junto a programación dinámica, producen unos resultados excelentes, 
  pero no siempre son la mejor solución.
\end{itemize}
Como conclusión final de este estudio, y de esta primera sección, se establece el algoritmo de la
programación dinámica como la mejor técnica tradicional para la resolución del Problema del Viajante.

En este capítulo se cuenta tal y tal.

El capítulo se estructura en $n$ apartados\ldots


\section{Estado del Arte}
\label{sec:estadoarte}

En el estado del arte se enumeran los trabajos más relevantes de otros grupos de investigación. A continuación se muestra un ejemplo del uso de viñetas que nos proporciona \texttt{itemize}:

\begin{itemize}
  \item En el trabajo .....
  \item En el siguiente trabajo.....
\end{itemize}

O citas en un párrafo real: Sin embargo, hay entornos acústicos donde las tasas de error conseguidas son todavía demasiado altas. En concreto, las aplicaciones en las que la captura de la señal de habla se hace usando micrófonos alejados del locutor (típicamente para distancias superiores a un metro) muestran una fuerte sensibilidad a los problemas de reverberación, ruido aditivo y baja relación señal a ruido (\cite{gelbart02},\cite{kochkin02}). En estos entornos, se ha propuesto el uso de arrays de micrófonos como un método para mejorar la calidad del habla capturada \cite{seltzer03}\cite{herbordt05}.

Existen múltiples formas de insertar figuras en Latex. A continuación, se muestra un ejemplo del uso de \texttt{figure}. Como se puede ver en la Figura \ref{fig:fig1} también se pueden poner referencias a las figuras por medio de \texttt{ref} y la etiqueta \texttt{label} de la figura en particular.

\begin{figure}[h] %el especificador [h] indica que ponga la figura aqui si es posible
  \centering \includegraphics[width=4.7in]{Figure1}
  % where an .eps filename suffix will be assumed under latex, and a .pdf suffix will be assumed for pdflatex
  \caption{Departamento de Electrónica.}
  \label{fig:fig1}
\end{figure}

Y ahora un ejemplo en el que ponemos el \texttt{caption} en el lateral:

\begin{SCfigure}
  \centering \includegraphics[width=0.5\textwidth]{Figure1}
  \caption{Departamento de Electrónica en el lateral.}
\end{SCfigure}



\section{Técnicas utilizadas}
\label{sec:tecnicas-utilizadas}

Aquí vamos a probar todos los niveles de sección disponibles, para evaluar la asignación de \texttt{tocdepth}...

Blah, blah, blah\ldots


\subsection{Subsección}
\label{sec:subseccion}


\subsubsection{Subsubsección}
\label{sec:subsubseccion}

\paragraph{Paragraph}
\label{sec:paragraph-1}


\subparagraph{Subparagraph}
\label{sec:subparagraph}



\section{Conclusiones}
\label{sec:conclusiones-teoria}

Blah, blah, blah\ldots

%%% Local Variables:
%%% TeX-master: "../book"
%%% End:
